%% 
%% Copyright 2007-2025 Elsevier Ltd
%% 
%% This file is part of the 'Elsarticle Bundle'.
%% ---------------------------------------------
%% 
%% It may be distributed under the conditions of the LaTeX Project Public
%% License, either version 1.3 of this license or (at your option) any
%% later version.  The latest version of this license is in
%%    http://www.latex-project.org/lppl.txt
%% and version 1.3 or later is part of all distributions of LaTeX
%% version 1999/12/01 or later.
%% 
%% The list of all files belonging to the 'Elsarticle Bundle' is
%% given in the file `manifest.txt'.
%% 
%% Template article for Elsevier's document class `elsarticle'
%% with numbered style bibliographic references
%% SP 2008/03/01
%% $Id: elsarticle-template-num.tex 272 2025-01-09 17:36:26Z rishi $
%%
\documentclass[preprint,12pt]{elsarticle}

%% Use the option review to obtain double line spacing
%% \documentclass[authoryear,preprint,review,12pt]{elsarticle}

%% Use the options 1p,twocolumn; 3p; 3p,twocolumn; 5p; or 5p,twocolumn
%% for a journal layout:
%% \documentclass[final,1p,times]{elsarticle}
%% \documentclass[final,1p,times,twocolumn]{elsarticle}
%% \documentclass[final,3p,times]{elsarticle}
%% \documentclass[final,3p,times,twocolumn]{elsarticle}
%% \documentclass[final,5p,times]{elsarticle}
%% \documentclass[final,5p,times,twocolumn]{elsarticle}

%% For including figures, graphicx.sty has been loaded in
%% elsarticle.cls. If you prefer to use the old commands
%% please give \usepackage{epsfig}

%% The amssymb package provides various useful mathematical symbols
\usepackage{amssymb}
%% The amsmath package provides various useful equation environments.
\usepackage{amsmath}
%% The amsthm package provides extended theorem environments
%% \usepackage{amsthm}

\usepackage{xcolor}

%% The lineno packages adds line numbers. Start line numbering with
%% \begin{linenumbers}, end it with \end{linenumbers}. Or switch it on
%% for the whole article with \linenumbers.
%% \usepackage{lineno}

\journal{Nuclear Physics B}

\begin{document}

\begin{frontmatter}

%% Title, authors and addresses

%% use the tnoteref command within \title for footnotes;
%% use the tnotetext command for theassociated footnote;
%% use the fnref command within \author or \affiliation for footnotes;
%% use the fntext command for theassociated footnote;
%% use the corref command within \author for corresponding author footnotes;
%% use the cortext command for theassociated footnote;
%% use the ead command for the email address,
%% and the form \ead[url] for the home page:
%% \title{Title\tnoteref{label1}}
%% \tnotetext[label1]{}
%% \author{Name\corref{cor1}\fnref{label2}}
%% \ead{email address}
%% \ead[url]{home page}
%% \fntext[label2]{}
%% \cortext[cor1]{}
%% \affiliation{organization={},
%%             addressline={},
%%             city={},
%%             postcode={},
%%             state={},
%%             country={}}
%% \fntext[label3]{}

\title{}

%% use optional labels to link authors explicitly to addresses:
%% \author[label1,label2]{}
%% \affiliation[label1]{organization={},
%%             addressline={},
%%             city={},
%%             postcode={},
%%             state={},
%%             country={}}
%%
%% \affiliation[label2]{organization={},
%%             addressline={},
%%             city={},
%%             postcode={},
%%             state={},
%%             country={}}

\author{Anatoly Belikov} %% Author name

%% Author affiliation
\affiliation{organization={SingularityNet},%Department and Organization
            addressline={}, 
            city={},
            postcode={}, 
            state={},
            country={}}

%% Abstract
\begin{abstract}
%% Text of abstract
Abstract text.
\end{abstract}

%%Graphical abstract
\begin{graphicalabstract}
%\includegraphics{grabs}
\end{graphicalabstract}

%%Research highlights
\begin{highlights}
\item Research highlight 1
\item Research highlight 2
\end{highlights}

%% Keywords
\begin{keyword}
%% keywords here, in the form: keyword \sep keyword

%% PACS codes here, in the form: \PACS code \sep code

%% MSC codes here, in the form: \MSC code \sep code
%% or \MSC[2008] code \sep code (2000 is the default)

\end{keyword}

\end{frontmatter}

%% Add \usepackage{lineno} before \begin{document} and uncomment 
%% following line to enable line numbers
%% \linenumbers

%% main text
%%

\maketitle

\section*{Reinforce}

\subsection*{Trajectory or path or episode}

$\tau = (s_0, a_0, r_0, s_1, a_1, r_1, \dots, s_T)$

Where:
\begin{itemize}
	\item $s_t$: state at time $t$
	\item $a_t$: action
	\item $r_t$: reward
	\item $T$: episode length
\end{itemize}

Policy: $\pi_\theta(a_t | s_t)$

\subsection*{Discounted return}

\textcolor{red}{todo - infinite vs finite horizon}

$R(\tau) = \sum_{t=0}^T \gamma^t r_t$

This is a discounted return over the whole trajectory. For infinite trajectory we should use $\gamma < 1$ to make it a convergent series, but in practice it’s possible to omit this constant e.g. when episodes are short or there is just one final reward etc..

Future discounted returns from time $t$:

$G_t = \gamma^t \sum_{k=t}^T \gamma^{k-t} r_k$ \quad or \quad $G_t = \sum_{k=t}^T \gamma^k r_k$

Objective (expected return):

$J(\theta) = \mathbb{E}_{\tau \sim \pi_\theta(\tau)} [R(\tau)]$

$J(\theta) = \int \pi_\theta(\tau)R(\tau) d\tau$

\section*{Policy gradient derivation}
Take gradient of the objective:
$\nabla_\theta J(\theta) = \int \nabla_\theta \pi_\theta(\tau)R(\tau) d\tau$

Log-derivative trick:

$log' f(x)= \frac {f'(x) } {f(x)} => f'(x) =f(x) log' f(x)$

%
Using this identity we get:

$\nabla_\theta J(\theta) = \int_{\tau \sim \pi_\theta(\tau)} 
\pi_\theta(\tau) \nabla_\theta \log \pi_\theta(\tau) R(\tau)$

Replace integral with expectation:

$\nabla_\theta J(\theta) = \mathbb{E}_{\tau \sim \pi_\theta(\tau)} [\nabla_\theta \log \pi_\theta(\tau) R(\tau)]$

In this equation we have for $\pi_\theta(\tau)$:

$\pi_\theta(\tau) = p(s_0) \prod_{t=0}^{T-1} \pi_\theta(a_t|s_t) p(s_{t+1}|s_t, a_t)$

if we take log of this product we get sum:

$\log \pi_\theta(\tau) = \log p(s_0) +  \sum_{t=0}^{T-1} \log \pi_\theta(a_t|s_t) + \sum_{t=0}^{T-1} \log  p(s_{t+1}|s_t, a_t)
$

Now take gradient with respect to $\theta$, terms not depending on $\theta$ drop.

$\nabla_\theta \log \pi_\theta(\tau) = \sum_{t=0}^T \nabla_\theta \log \pi_\theta(a_t | s_t)$

Plug this expression back to the equation:

$\nabla_\theta J(\theta) = \mathbb{E}_{\tau \sim \pi_\theta(\tau)}  \sum_{t=0}^T \nabla_\theta \log \pi_\theta(a_t | s_t)  R(\tau) $


Since past rewards don’t depend on current(or future) action at  we can replace $R(\tau)$ with $G_t$.

\textcolor{red}{proof - todo} 
 
$\nabla_\theta J(\theta) = \mathbb{E}_{\tau \sim \pi_\theta(\tau)}  \sum_{t=0}^T \nabla_\theta \log \pi_\theta(a_t | s_t)  G_t $

\section{Q and Value Functions}
Now we can extract Q and value functions from the expection above:

$Q_{\pi_\theta}(s_t, a_t) = \mathbb{E}_{\tau \sim \pi_\theta(\tau)}  \sum_{k=t}^T \gamma^{k-t} r_k | s_t, a_t $

$V_{\pi_\theta}(s_t) = \mathbb{E}_{a \sim \pi_\theta(s)} Q_{\pi_\theta}(a_t, s_t)$

This is a relative time which allows for bellman expectation form(see below in Useful identities section).

With $G_t = \gamma^t \sum_{k=t}^T \gamma^{k-t} r_k$

we can rewrite out gradient to 


$\nabla_\theta J(\theta) = \mathbb{E}_{\tau \sim \pi_\theta(\tau)}  \sum_{t=0}^T \gamma^t \nabla_\theta \log \pi_\theta(a_t | s_t) Q^{\pi_\theta}(s_t, a_t) $


$\gamma^t$ comes from $G_t$

\textcolor{blue}{
in practice we can drop $\gamma^t$ or use different future discount functions e.g $w(t) = 1/(1 + t/d)$.}


$\gamma^t$ might drop to too small values for long episodes. For example:


$\gamma=0.95^{t=150} \approx 1500.00045$ vs

$\frac{1} {1+ \frac{t=150}{d=2}} \approx  0.013$ 


\section*{State distribution view}
We can marginalise distribution of trajectories $\pi_\theta(\tau)$ to $d^{\pi_\theta}_t(s)$ - distribution of states at time $t$ for our policy.

Factor one term for time $t$ inside trajectory summation:

$\nabla_\theta \log \pi_\theta(a_t|s_t) Q^{\pi_\theta}(s_t,a_t)$

becomes:

$\sum_s \sum_a P(s_t=s, a_t=a) \nabla_\theta \log \pi_\theta(a|s) Q^{\pi_\theta}(s,a)$

Given that

$P(s_t=s, a_t=a) = P(s_t=s)\pi_\theta(a|s)$,

we get

$\sum_s P(s_t=s) \sum_a \pi_\theta(a|s) \nabla_\theta \log \pi_\theta(a|s) Q^{\pi_\theta}(s,a)$

This is where the state distribution appears.

Put the factored equation into $\nabla_\theta J(\theta)$:

$ \sum_s \left( \sum_{t=0}^T \gamma^t P(s_t=s) \right) \sum_a \pi_\theta(a|s) \nabla_\theta \log \pi_\theta(a|s) Q^{\pi_\theta}(s,a)$

Define the discounted state visitation distribution:

$d_\gamma^{\pi_\theta}(s) = \sum_{t=0}^T \gamma^t P(s_t=s)$

Note: it is unnormalised.

Proper distribution:

$d^{\pi_\theta}(s) = (1-\gamma) \sum_{t=0}^{T} \gamma^t d^{\pi_\theta}_t(s)$

Then

$\nabla_\theta J(\theta) = \sum_s d^{\pi_\theta}(s) \sum_a \pi_\theta(a|s) \nabla_\theta \log \pi_\theta(a|s) Q^{\pi_\theta}(s,a)$

Given that

$\pi_\theta(a|s) \nabla_\theta \log \pi_\theta(a|s) = \nabla_\theta \pi_\theta(a|s)$,

we can also write

$\nabla_\theta J(\theta) = \sum_s d^{\pi_\theta}(s) \sum_a \nabla_\theta \pi_\theta(a|s) Q^{\pi_\theta}(s,a)$

We can use the formulation above and explicitly sum over actions if possible, but more often in practice we sample actions, so the $\log \pi_\theta$ form comes back:

$\nabla_\theta J(\theta) = \mathbb{E}_{s \sim d^{\pi_\theta}} \mathbb{E}_{a \sim \pi_\theta(\cdot|s)} [\nabla_\theta \log \pi_\theta(a|s) Q^{\pi_\theta}(s,a)]$

This is called the policy gradient theorem.

Effectively we can replace expectations over trajectories with expectations over state distribution under current policy.


\section{Generalized Advantage Estimation (GAE-$\lambda$) }

\textcolor{red}{todo}


\section*{Useful identities}


Bellman expectation form:


$V^{\pi_\theta}(s) = \mathbb{E}_{\tau \sim \pi_\theta(\tau)} [r(s,a) + \gamma V^\pi(s')]$

$Q^{\pi_\theta}(s,a) = \mathbb{E}_{\tau \sim \pi_\theta(\tau)} [r_t + \gamma V^{\pi_\theta}(s_{t+1}) | s_t=s, a_t=a]$

$Q^{\pi_\theta}(s,a) = \mathbb{E}_{\tau \sim \pi_\theta(\tau)} [r_t + \gamma \mathbb{E}_{a' \sim \pi_\theta(\cdot|s')} [Q^{\pi_\theta}(s', a')]]$


With relative time $G_t$:

$G_t = r_t + \gamma r_{t+1} + \gamma ^2 r_{t+2} + ...$

$V^{\pi_\theta}(s) = \mathbb{E}_{\pi_\theta} \left[ G_t|s_t=s \right]$

\end{document}

\endinput
%%
%% End of file `elsarticle-template-num.tex'.
